\scanpage{78}
\unsure{essentielle triviale|essentielles triviales}, qui fait penser à (6), est
\[
\tag{28}{}^{uv}l={}^u({}^vl)
\qquad u,v\in\Gamma,\quad l:x\longrightarrow y\quad\text{chemin dans }\widehat\Pi_1(X),
\]
\deleted{d'où,} si $l$ est chemin de $\Pi_1(X)$, d'où
\[
{}^vl=g_l(v)\cdot l
\]
\deleted{\unsure{grâce à|qui est} (6)} d'où on tire par (6)
\[
{}^u({}^vl)={}^u\{g_l(v)\}\cdot\underbrace{{}^ul}_{g_l(u)\cdot l}
={}^ug_l(v)g_l(u)\cdot l
\]
d'où on tire
\[
\tag{29}\left\{\begin{aligned}
\deleted{\unsure{u}}g_l(uv)&={}^ug_l(v)\cdot g_l(u)\qquad\text{soit}\\
g_l^{-1}(uv)&=g_l^{-1}(u)\cdot{}^ug_l^{-1}(v)
\end{aligned}\right.
\]
\unsure{Et de $=$}
\[
\tag{30}f_l(uv)=f_l(u)\cdot{}^uf_l(v)
\]
i.e. $\unsure{f_{l\alpha}|f_l}$ et $g_l^{-1}$ sont, pour $l$ fixé, $\unsure{u|\text{en}}$ variable, des $1$-cocycles \inserted{sur $\Gamma$} à valeurs dans $\widehat\Pi_1(X,x)$ resp. $\widehat\Pi_1(X,y)$, \inserted{(groupes sur lesquels $\Gamma$ opère)}, d'ailleurs déduits l'un de l'autre par l'isomorphisme
\diagram{\begin{tikzcd}[ampersand replacement=\&]
\widehat\Pi_1(X,x)\arrow[r,no head,"\Pi_1(l)"]\&\widehat\Pi_1(X,y)
\end{tikzcd}}
induit par $l$ \deleted{\unsure{en $l$}}. Ici, c'est la formule pour $f_l$ qui a une gueule plus sympathique que celle pour $g_l$ (pour (10)(11) c'était l'inverse --- \inserted{exprim.} ``multiplicativité en $l$''). Les $1$-cocycles obtenus (\deleted{et} $f$ \inserted{un $1$-\unsure{cocycles|cocycle}} \unsure{droits|droit}, $g$ un $1$-cocycle à gauche) correspondent à des $\Gamma$ torseurs à droite resp. à gauche sous $\widehat\Pi_1(X,x)$ resp! $\widehat\Pi_1(X,y)$, qui bien sûr n'est autre que $\operatorname{Hom}_{\widehat\Pi_1(X)}(x,y)$

\margin{On pourrait expliciter \unsure{un de vue qui voit le|au demeurant le} couple $(f_l,g_l)$ comme un ``bicocycle'' sur $\Gamma$, relatif à un $\Gamma$-bitorseur $\operatorname{Hom}_{\widehat\Pi_1(X)}(x,y)$ sous $(\widehat\Pi_1(X,y),\widehat\Pi_1(X,x))$.}
\margin{77}
